% Options for packages loaded elsewhere
\PassOptionsToPackage{unicode}{hyperref}
\PassOptionsToPackage{hyphens}{url}
%
\documentclass[
]{article}
\usepackage{amsmath,amssymb}
\usepackage{iftex}
\ifPDFTeX
  \usepackage[T1]{fontenc}
  \usepackage[utf8]{inputenc}
  \usepackage{textcomp} % provide euro and other symbols
\else % if luatex or xetex
  \usepackage{unicode-math} % this also loads fontspec
  \defaultfontfeatures{Scale=MatchLowercase}
  \defaultfontfeatures[\rmfamily]{Ligatures=TeX,Scale=1}
\fi
\usepackage{lmodern}
\ifPDFTeX\else
  % xetex/luatex font selection
\fi
% Use upquote if available, for straight quotes in verbatim environments
\IfFileExists{upquote.sty}{\usepackage{upquote}}{}
\IfFileExists{microtype.sty}{% use microtype if available
  \usepackage[]{microtype}
  \UseMicrotypeSet[protrusion]{basicmath} % disable protrusion for tt fonts
}{}
\makeatletter
\@ifundefined{KOMAClassName}{% if non-KOMA class
  \IfFileExists{parskip.sty}{%
    \usepackage{parskip}
  }{% else
    \setlength{\parindent}{0pt}
    \setlength{\parskip}{6pt plus 2pt minus 1pt}}
}{% if KOMA class
  \KOMAoptions{parskip=half}}
\makeatother
\usepackage{xcolor}
\usepackage[margin=1in]{geometry}
\usepackage{color}
\usepackage{fancyvrb}
\newcommand{\VerbBar}{|}
\newcommand{\VERB}{\Verb[commandchars=\\\{\}]}
\DefineVerbatimEnvironment{Highlighting}{Verbatim}{commandchars=\\\{\}}
% Add ',fontsize=\small' for more characters per line
\usepackage{framed}
\definecolor{shadecolor}{RGB}{248,248,248}
\newenvironment{Shaded}{\begin{snugshade}}{\end{snugshade}}
\newcommand{\AlertTok}[1]{\textcolor[rgb]{0.94,0.16,0.16}{#1}}
\newcommand{\AnnotationTok}[1]{\textcolor[rgb]{0.56,0.35,0.01}{\textbf{\textit{#1}}}}
\newcommand{\AttributeTok}[1]{\textcolor[rgb]{0.13,0.29,0.53}{#1}}
\newcommand{\BaseNTok}[1]{\textcolor[rgb]{0.00,0.00,0.81}{#1}}
\newcommand{\BuiltInTok}[1]{#1}
\newcommand{\CharTok}[1]{\textcolor[rgb]{0.31,0.60,0.02}{#1}}
\newcommand{\CommentTok}[1]{\textcolor[rgb]{0.56,0.35,0.01}{\textit{#1}}}
\newcommand{\CommentVarTok}[1]{\textcolor[rgb]{0.56,0.35,0.01}{\textbf{\textit{#1}}}}
\newcommand{\ConstantTok}[1]{\textcolor[rgb]{0.56,0.35,0.01}{#1}}
\newcommand{\ControlFlowTok}[1]{\textcolor[rgb]{0.13,0.29,0.53}{\textbf{#1}}}
\newcommand{\DataTypeTok}[1]{\textcolor[rgb]{0.13,0.29,0.53}{#1}}
\newcommand{\DecValTok}[1]{\textcolor[rgb]{0.00,0.00,0.81}{#1}}
\newcommand{\DocumentationTok}[1]{\textcolor[rgb]{0.56,0.35,0.01}{\textbf{\textit{#1}}}}
\newcommand{\ErrorTok}[1]{\textcolor[rgb]{0.64,0.00,0.00}{\textbf{#1}}}
\newcommand{\ExtensionTok}[1]{#1}
\newcommand{\FloatTok}[1]{\textcolor[rgb]{0.00,0.00,0.81}{#1}}
\newcommand{\FunctionTok}[1]{\textcolor[rgb]{0.13,0.29,0.53}{\textbf{#1}}}
\newcommand{\ImportTok}[1]{#1}
\newcommand{\InformationTok}[1]{\textcolor[rgb]{0.56,0.35,0.01}{\textbf{\textit{#1}}}}
\newcommand{\KeywordTok}[1]{\textcolor[rgb]{0.13,0.29,0.53}{\textbf{#1}}}
\newcommand{\NormalTok}[1]{#1}
\newcommand{\OperatorTok}[1]{\textcolor[rgb]{0.81,0.36,0.00}{\textbf{#1}}}
\newcommand{\OtherTok}[1]{\textcolor[rgb]{0.56,0.35,0.01}{#1}}
\newcommand{\PreprocessorTok}[1]{\textcolor[rgb]{0.56,0.35,0.01}{\textit{#1}}}
\newcommand{\RegionMarkerTok}[1]{#1}
\newcommand{\SpecialCharTok}[1]{\textcolor[rgb]{0.81,0.36,0.00}{\textbf{#1}}}
\newcommand{\SpecialStringTok}[1]{\textcolor[rgb]{0.31,0.60,0.02}{#1}}
\newcommand{\StringTok}[1]{\textcolor[rgb]{0.31,0.60,0.02}{#1}}
\newcommand{\VariableTok}[1]{\textcolor[rgb]{0.00,0.00,0.00}{#1}}
\newcommand{\VerbatimStringTok}[1]{\textcolor[rgb]{0.31,0.60,0.02}{#1}}
\newcommand{\WarningTok}[1]{\textcolor[rgb]{0.56,0.35,0.01}{\textbf{\textit{#1}}}}
\usepackage{graphicx}
\makeatletter
\def\maxwidth{\ifdim\Gin@nat@width>\linewidth\linewidth\else\Gin@nat@width\fi}
\def\maxheight{\ifdim\Gin@nat@height>\textheight\textheight\else\Gin@nat@height\fi}
\makeatother
% Scale images if necessary, so that they will not overflow the page
% margins by default, and it is still possible to overwrite the defaults
% using explicit options in \includegraphics[width, height, ...]{}
\setkeys{Gin}{width=\maxwidth,height=\maxheight,keepaspectratio}
% Set default figure placement to htbp
\makeatletter
\def\fps@figure{htbp}
\makeatother
\setlength{\emergencystretch}{3em} % prevent overfull lines
\providecommand{\tightlist}{%
  \setlength{\itemsep}{0pt}\setlength{\parskip}{0pt}}
\setcounter{secnumdepth}{-\maxdimen} % remove section numbering
\ifLuaTeX
  \usepackage{selnolig}  % disable illegal ligatures
\fi
\usepackage{bookmark}
\IfFileExists{xurl.sty}{\usepackage{xurl}}{} % add URL line breaks if available
\urlstyle{same}
\hypersetup{
  pdftitle={Ind 4: Методи ресемплінгу},
  hidelinks,
  pdfcreator={LaTeX via pandoc}}

\title{Ind 4: Методи ресемплінгу}
\usepackage{etoolbox}
\makeatletter
\providecommand{\subtitle}[1]{% add subtitle to \maketitle
  \apptocmd{\@title}{\par {\large #1 \par}}{}{}
}
\makeatother
\subtitle{Виконала: Єршова Ганна, студентка групи ПМіМ-13}
\author{}
\date{\vspace{-2.5em}}

\begin{document}
\maketitle

Встановлюємо значення variant, seed та redundant

\begin{Shaded}
\begin{Highlighting}[]
\FunctionTok{library}\NormalTok{(boot)}
\NormalTok{group }\OtherTok{\textless{}{-}} \DecValTok{3}
\NormalTok{variant }\OtherTok{\textless{}{-}}\NormalTok{ group }\SpecialCharTok{*} \DecValTok{25} \SpecialCharTok{+} \DecValTok{4}
\FunctionTok{set.seed}\NormalTok{(variant)}
\NormalTok{redundant }\OtherTok{\textless{}{-}} \FunctionTok{round}\NormalTok{(}\FunctionTok{runif}\NormalTok{(}\DecValTok{1}\NormalTok{, }\DecValTok{5} \SpecialCharTok{+}\NormalTok{ group, }\DecValTok{25} \SpecialCharTok{{-}}\NormalTok{ group))}
\end{Highlighting}
\end{Shaded}

\subsection{Завдання
№1}\label{ux437ux430ux432ux434ux430ux43dux43dux44f-1}

Завантажуємо дані та видаляємо redundant \% значень

\begin{Shaded}
\begin{Highlighting}[]
\FunctionTok{library}\NormalTok{(MASS)}
\NormalTok{boston\_new }\OtherTok{\textless{}{-}}\NormalTok{ Boston[}\SpecialCharTok{{-}}\FunctionTok{sample}\NormalTok{(}\DecValTok{1}\SpecialCharTok{:}\FunctionTok{length}\NormalTok{(Boston[,}\DecValTok{1}\NormalTok{]), }\FunctionTok{round}\NormalTok{((redundant}\SpecialCharTok{/}\DecValTok{100}\NormalTok{) }\SpecialCharTok{*} \FunctionTok{length}\NormalTok{(Boston[,}\DecValTok{1}\NormalTok{]))),]}
\end{Highlighting}
\end{Shaded}

\begin{Shaded}
\begin{Highlighting}[]
\NormalTok{?Boston}
\end{Highlighting}
\end{Shaded}

Додамо змінну, яка позначатиме, чи рівень злочинності перевищує середній

\begin{Shaded}
\begin{Highlighting}[]
\NormalTok{crime\_mean }\OtherTok{\textless{}{-}} \FunctionTok{mean}\NormalTok{(boston\_new}\SpecialCharTok{$}\NormalTok{crim)}
\NormalTok{boston\_new}\SpecialCharTok{$}\NormalTok{crime\_level }\OtherTok{\textless{}{-}} \FunctionTok{ifelse}\NormalTok{(boston\_new}\SpecialCharTok{$}\NormalTok{crim }\SpecialCharTok{\textgreater{}}\NormalTok{ crime\_mean, }\DecValTok{1}\NormalTok{, }\DecValTok{0}\NormalTok{)}
\end{Highlighting}
\end{Shaded}

Будуємо модель логістичної регресії для передбачення рівня злочинності
більшого чи меншого за середній на основі змінних nox та rad

\begin{Shaded}
\begin{Highlighting}[]
\NormalTok{train }\OtherTok{\textless{}{-}} \FunctionTok{sample}\NormalTok{(}\FunctionTok{nrow}\NormalTok{(boston\_new), }\FunctionTok{nrow}\NormalTok{(boston\_new)}\SpecialCharTok{/}\DecValTok{2}\NormalTok{)}
\NormalTok{lr.train.fit }\OtherTok{\textless{}{-}} \FunctionTok{glm}\NormalTok{(crime\_level }\SpecialCharTok{\textasciitilde{}}\NormalTok{ nox }\SpecialCharTok{+}\NormalTok{ rad, }\AttributeTok{data=}\NormalTok{boston\_new, }\AttributeTok{family=}\NormalTok{binomial, }\AttributeTok{subset=}\NormalTok{train)}
\FunctionTok{summary}\NormalTok{(lr.train.fit)}
\end{Highlighting}
\end{Shaded}

\begin{verbatim}
## 
## Call:
## glm(formula = crime_level ~ nox + rad, family = binomial, data = boston_new, 
##     subset = train)
## 
## Coefficients:
##             Estimate Std. Error z value Pr(>|z|)    
## (Intercept) -19.7546     9.2215  -2.142 0.032175 *  
## nox          17.6187     9.9829   1.765 0.077583 .  
## rad           0.4782     0.1422   3.363 0.000771 ***
## ---
## Signif. codes:  0 '***' 0.001 '**' 0.01 '*' 0.05 '.' 0.1 ' ' 1
## 
## (Dispersion parameter for binomial family taken to be 1)
## 
##     Null deviance: 201.049  on 201  degrees of freedom
## Residual deviance:  20.897  on 199  degrees of freedom
## AIC: 26.897
## 
## Number of Fisher Scoring iterations: 10
\end{verbatim}

\begin{Shaded}
\begin{Highlighting}[]
\NormalTok{lr.prob }\OtherTok{\textless{}{-}} \FunctionTok{predict}\NormalTok{(lr.train.fit, boston\_new)}
\NormalTok{lr.prob.val }\OtherTok{\textless{}{-}}\NormalTok{ lr.prob[}\SpecialCharTok{{-}}\NormalTok{train]}
\NormalTok{lr.pred }\OtherTok{\textless{}{-}} \FunctionTok{rep}\NormalTok{(}\DecValTok{0}\NormalTok{, }\FunctionTok{length}\NormalTok{(lr.prob.val))}
\NormalTok{lr.pred[lr.prob.val }\SpecialCharTok{\textgreater{}} \FloatTok{0.5}\NormalTok{] }\OtherTok{\textless{}{-}} \DecValTok{1}
\end{Highlighting}
\end{Shaded}

Тестова помилка

\begin{Shaded}
\begin{Highlighting}[]
\FunctionTok{mean}\NormalTok{((boston\_new}\SpecialCharTok{$}\NormalTok{crime\_level[}\SpecialCharTok{{-}}\NormalTok{train] }\SpecialCharTok{!=}\NormalTok{ lr.pred))}
\end{Highlighting}
\end{Shaded}

\begin{verbatim}
## [1] 0.01477833
\end{verbatim}

Повторимо попередню процедуру три рази

\begin{Shaded}
\begin{Highlighting}[]
\NormalTok{train.model }\OtherTok{\textless{}{-}} \ControlFlowTok{function}\NormalTok{(seed) \{}
  \FunctionTok{set.seed}\NormalTok{(seed)}
\NormalTok{  t.split }\OtherTok{\textless{}{-}} \FunctionTok{sample}\NormalTok{(}\FunctionTok{nrow}\NormalTok{(boston\_new), }\FunctionTok{nrow}\NormalTok{(boston\_new) }\SpecialCharTok{*} \FloatTok{0.5}\NormalTok{)}
\NormalTok{  lr }\OtherTok{\textless{}{-}} \FunctionTok{glm}\NormalTok{(crime\_level }\SpecialCharTok{\textasciitilde{}}\NormalTok{ nox }\SpecialCharTok{+}\NormalTok{ rad, }\AttributeTok{data =}\NormalTok{ boston\_new, }\AttributeTok{family =}\NormalTok{ binomial, }\AttributeTok{subset =}\NormalTok{ t.split)}
\NormalTok{  lr.p }\OtherTok{\textless{}{-}} \FunctionTok{predict}\NormalTok{(lr, boston\_new[}\SpecialCharTok{{-}}\NormalTok{t.split, ], }\AttributeTok{type =} \StringTok{"response"}\NormalTok{)}
\NormalTok{  pred }\OtherTok{\textless{}{-}} \FunctionTok{ifelse}\NormalTok{(lr.p }\SpecialCharTok{\textgreater{}} \FloatTok{0.5}\NormalTok{, }\DecValTok{1}\NormalTok{, }\DecValTok{0}\NormalTok{)}
\NormalTok{  test\_error }\OtherTok{\textless{}{-}} \FunctionTok{mean}\NormalTok{(boston\_new}\SpecialCharTok{$}\NormalTok{crime\_level[}\SpecialCharTok{{-}}\NormalTok{t.split] }\SpecialCharTok{!=}\NormalTok{ pred)}
  \FunctionTok{return}\NormalTok{(test\_error)}
\NormalTok{\}}
\end{Highlighting}
\end{Shaded}

\begin{Shaded}
\begin{Highlighting}[]
\FunctionTok{set.seed}\NormalTok{(variant)}
\NormalTok{seeds }\OtherTok{\textless{}{-}} \FunctionTok{sample}\NormalTok{(}\DecValTok{1}\SpecialCharTok{:}\DecValTok{1000}\NormalTok{, }\DecValTok{3}\NormalTok{)}
\NormalTok{test\_errors }\OtherTok{\textless{}{-}} \FunctionTok{sapply}\NormalTok{(seeds, train.model)}
\end{Highlighting}
\end{Shaded}

\begin{verbatim}
## Warning: glm.fit: fitted probabilities numerically 0 or 1 occurred
\end{verbatim}

\begin{Shaded}
\begin{Highlighting}[]
\FunctionTok{print}\NormalTok{(test\_errors)}
\end{Highlighting}
\end{Shaded}

\begin{verbatim}
## [1] 0.01970443 0.02463054 0.01477833
\end{verbatim}

\begin{Shaded}
\begin{Highlighting}[]
\NormalTok{average\_test\_error }\OtherTok{\textless{}{-}} \FunctionTok{mean}\NormalTok{(test\_errors)}
\FunctionTok{cat}\NormalTok{(}\StringTok{"Average test error over three runs:"}\NormalTok{, average\_test\_error, }\StringTok{"}\SpecialCharTok{\textbackslash{}n}\StringTok{"}\NormalTok{)}
\end{Highlighting}
\end{Shaded}

\begin{verbatim}
## Average test error over three runs: 0.01970443
\end{verbatim}

Додаємо змінну medv

\begin{Shaded}
\begin{Highlighting}[]
\NormalTok{train.model.medv }\OtherTok{\textless{}{-}} \ControlFlowTok{function}\NormalTok{(seed) \{}
  \FunctionTok{set.seed}\NormalTok{(seed)}
\NormalTok{  t.split }\OtherTok{\textless{}{-}} \FunctionTok{sample}\NormalTok{(}\FunctionTok{nrow}\NormalTok{(boston\_new), }\FunctionTok{nrow}\NormalTok{(boston\_new) }\SpecialCharTok{*} \FloatTok{0.5}\NormalTok{)}
\NormalTok{  lr }\OtherTok{\textless{}{-}} \FunctionTok{glm}\NormalTok{(crime\_level }\SpecialCharTok{\textasciitilde{}}\NormalTok{ nox }\SpecialCharTok{+}\NormalTok{ rad }\SpecialCharTok{+}\NormalTok{ medv, }\AttributeTok{data =}\NormalTok{ boston\_new, }\AttributeTok{family =}\NormalTok{ binomial, }\AttributeTok{subset =}\NormalTok{ t.split)}
\NormalTok{  lr.p }\OtherTok{\textless{}{-}} \FunctionTok{predict}\NormalTok{(lr, boston\_new[}\SpecialCharTok{{-}}\NormalTok{t.split, ], }\AttributeTok{type =} \StringTok{"response"}\NormalTok{)}
\NormalTok{  pred }\OtherTok{\textless{}{-}} \FunctionTok{ifelse}\NormalTok{(lr.p }\SpecialCharTok{\textgreater{}} \FloatTok{0.5}\NormalTok{, }\DecValTok{1}\NormalTok{, }\DecValTok{0}\NormalTok{)}
\NormalTok{  test\_error }\OtherTok{\textless{}{-}} \FunctionTok{mean}\NormalTok{(boston\_new}\SpecialCharTok{$}\NormalTok{crime\_level[}\SpecialCharTok{{-}}\NormalTok{t.split] }\SpecialCharTok{!=}\NormalTok{ pred)}
  \FunctionTok{return}\NormalTok{(test\_error)}
\NormalTok{\}}
\end{Highlighting}
\end{Shaded}

\begin{Shaded}
\begin{Highlighting}[]
\NormalTok{test\_errors }\OtherTok{\textless{}{-}} \FunctionTok{sapply}\NormalTok{(seeds, train.model.medv)}
\end{Highlighting}
\end{Shaded}

\begin{verbatim}
## Warning: glm.fit: fitted probabilities numerically 0 or 1 occurred
\end{verbatim}

\begin{Shaded}
\begin{Highlighting}[]
\FunctionTok{print}\NormalTok{(test\_errors)}
\end{Highlighting}
\end{Shaded}

\begin{verbatim}
## [1] 0.01970443 0.02463054 0.01477833
\end{verbatim}

\begin{Shaded}
\begin{Highlighting}[]
\NormalTok{average\_test\_error }\OtherTok{\textless{}{-}} \FunctionTok{mean}\NormalTok{(test\_errors)}
\FunctionTok{cat}\NormalTok{(}\StringTok{"Average test error over three runs:"}\NormalTok{, average\_test\_error, }\StringTok{"}\SpecialCharTok{\textbackslash{}n}\StringTok{"}\NormalTok{)}
\end{Highlighting}
\end{Shaded}

\begin{verbatim}
## Average test error over three runs: 0.01970443
\end{verbatim}

\subsection{Завдання
№2}\label{ux437ux430ux432ux434ux430ux43dux43dux44f-2}

Завантажуємо дані

\begin{Shaded}
\begin{Highlighting}[]
\NormalTok{auto }\OtherTok{\textless{}{-}} \FunctionTok{read.csv}\NormalTok{(}\StringTok{"Auto.csv"}\NormalTok{, }\AttributeTok{stringsAsFactors =} \ConstantTok{TRUE}\NormalTok{)}
\NormalTok{auto\_new }\OtherTok{\textless{}{-}}\NormalTok{ auto[}\SpecialCharTok{{-}}\FunctionTok{sample}\NormalTok{(}\DecValTok{1}\SpecialCharTok{:}\FunctionTok{length}\NormalTok{(auto[,}\DecValTok{1}\NormalTok{]), }\FunctionTok{round}\NormalTok{((redundant}\SpecialCharTok{/}\DecValTok{100}\NormalTok{) }\SpecialCharTok{*} \FunctionTok{length}\NormalTok{(auto[,}\DecValTok{1}\NormalTok{]))),]}
\end{Highlighting}
\end{Shaded}

Оцінка середнього значення та стандартної похибки для змінної mpg

\begin{Shaded}
\begin{Highlighting}[]
\NormalTok{mpg\_mean }\OtherTok{\textless{}{-}} \FunctionTok{mean}\NormalTok{(auto\_new}\SpecialCharTok{$}\NormalTok{mpg)}
\NormalTok{mpg\_se }\OtherTok{\textless{}{-}} \FunctionTok{sd}\NormalTok{(auto\_new}\SpecialCharTok{$}\NormalTok{mpg) }\SpecialCharTok{/} \FunctionTok{sqrt}\NormalTok{(}\FunctionTok{length}\NormalTok{(auto\_new}\SpecialCharTok{$}\NormalTok{mpg))}

\FunctionTok{cat}\NormalTok{(}\StringTok{"Середнє mpg:"}\NormalTok{, mpg\_mean, }\StringTok{"}\SpecialCharTok{\textbackslash{}n}\StringTok{"}\NormalTok{)}
\end{Highlighting}
\end{Shaded}

\begin{verbatim}
## Середнє mpg: 23.57233
\end{verbatim}

\begin{Shaded}
\begin{Highlighting}[]
\FunctionTok{cat}\NormalTok{(}\StringTok{"Стандартна похибка:"}\NormalTok{, mpg\_se, }\StringTok{"}\SpecialCharTok{\textbackslash{}n}\StringTok{"}\NormalTok{)}
\end{Highlighting}
\end{Shaded}

\begin{verbatim}
## Стандартна похибка: 0.4369786
\end{verbatim}

Оцінюємо стандартну похибку за допомогою бутстрапу

\begin{Shaded}
\begin{Highlighting}[]
\NormalTok{boot\_mean }\OtherTok{\textless{}{-}} \ControlFlowTok{function}\NormalTok{(data, indices) }\FunctionTok{mean}\NormalTok{(data[indices])}
\NormalTok{mpg\_sd\_boot }\OtherTok{\textless{}{-}} \FunctionTok{boot}\NormalTok{(}\AttributeTok{data =}\NormalTok{ auto\_new}\SpecialCharTok{$}\NormalTok{mpg, }\AttributeTok{statistic =}\NormalTok{ boot\_mean, }\AttributeTok{R =} \DecValTok{1000}\NormalTok{)}
\FunctionTok{cat}\NormalTok{(}\StringTok{"Бутстрап{-}оцінка стандартної похибки:"}\NormalTok{, }\FunctionTok{sd}\NormalTok{(mpg\_sd\_boot}\SpecialCharTok{$}\NormalTok{t), }\StringTok{"}\SpecialCharTok{\textbackslash{}n}\StringTok{"}\NormalTok{)}
\end{Highlighting}
\end{Shaded}

\begin{verbatim}
## Бутстрап-оцінка стандартної похибки: 0.4297621
\end{verbatim}

Оцінюємо стандартну похибку медіани

\begin{Shaded}
\begin{Highlighting}[]
\NormalTok{boot\_median }\OtherTok{\textless{}{-}} \ControlFlowTok{function}\NormalTok{(data, index) }\FunctionTok{median}\NormalTok{(data[index])}

\NormalTok{mpg\_median\_boot }\OtherTok{\textless{}{-}} \FunctionTok{boot}\NormalTok{(}\AttributeTok{data =}\NormalTok{ auto\_new}\SpecialCharTok{$}\NormalTok{mpg, }\AttributeTok{statistic =}\NormalTok{ boot\_median, }\AttributeTok{R =} \DecValTok{1000}\NormalTok{)}
\NormalTok{mpg\_median }\OtherTok{\textless{}{-}} \FunctionTok{median}\NormalTok{(auto\_new}\SpecialCharTok{$}\NormalTok{mpg)}
\FunctionTok{cat}\NormalTok{(}\StringTok{"Медіана mpg:"}\NormalTok{, mpg\_median, }\StringTok{"}\SpecialCharTok{\textbackslash{}n}\StringTok{"}\NormalTok{)}
\end{Highlighting}
\end{Shaded}

\begin{verbatim}
## Медіана mpg: 23
\end{verbatim}

\begin{Shaded}
\begin{Highlighting}[]
\FunctionTok{cat}\NormalTok{(}\StringTok{"Бутстрап{-}оцінка стандартної похибки медіани:"}\NormalTok{, }\FunctionTok{sd}\NormalTok{(mpg\_median\_boot}\SpecialCharTok{$}\NormalTok{t), }\StringTok{"}\SpecialCharTok{\textbackslash{}n}\StringTok{"}\NormalTok{)}
\end{Highlighting}
\end{Shaded}

\begin{verbatim}
## Бутстрап-оцінка стандартної похибки медіани: 0.8790468
\end{verbatim}

Оцінюємо стандартну похибку десятого процентиля

\begin{Shaded}
\begin{Highlighting}[]
\NormalTok{boot\_10 }\OtherTok{\textless{}{-}} \ControlFlowTok{function}\NormalTok{(data, index) }\FunctionTok{quantile}\NormalTok{(data[index], }\FloatTok{0.1}\NormalTok{)}

\NormalTok{mpg\_10\_boot }\OtherTok{\textless{}{-}} \FunctionTok{boot}\NormalTok{(}\AttributeTok{data =}\NormalTok{ auto\_new}\SpecialCharTok{$}\NormalTok{mpg, }\AttributeTok{statistic =}\NormalTok{ boot\_10, }\AttributeTok{R =} \DecValTok{1000}\NormalTok{)}
\NormalTok{mpg\_10 }\OtherTok{\textless{}{-}} \FunctionTok{quantile}\NormalTok{(auto\_new}\SpecialCharTok{$}\NormalTok{mpg, }\FloatTok{0.1}\NormalTok{)}
\FunctionTok{cat}\NormalTok{(}\StringTok{"Десятий процентиль mpg:"}\NormalTok{, mpg\_10, }\StringTok{"}\SpecialCharTok{\textbackslash{}n}\StringTok{"}\NormalTok{)}
\end{Highlighting}
\end{Shaded}

\begin{verbatim}
## Десятий процентиль mpg: 14
\end{verbatim}

\begin{Shaded}
\begin{Highlighting}[]
\FunctionTok{cat}\NormalTok{(}\StringTok{"Бутстрап{-}оцінка стандартної похибки десятого процентиля:"}\NormalTok{, }\FunctionTok{sd}\NormalTok{(mpg\_10\_boot}\SpecialCharTok{$}\NormalTok{t), }\StringTok{"}\SpecialCharTok{\textbackslash{}n}\StringTok{"}\NormalTok{)}
\end{Highlighting}
\end{Shaded}

\begin{verbatim}
## Бутстрап-оцінка стандартної похибки десятого процентиля: 0.4611278
\end{verbatim}

\subsection{Завдання
№3}\label{ux437ux430ux432ux434ux430ux43dux43dux44f-3}

Генеруємо дані

\begin{Shaded}
\begin{Highlighting}[]
\FunctionTok{set.seed}\NormalTok{(variant)}
\NormalTok{x }\OtherTok{=} \FunctionTok{rnorm}\NormalTok{ (}\DecValTok{100}\NormalTok{)}
\NormalTok{y }\OtherTok{=}\NormalTok{ variant}\SpecialCharTok{*}\NormalTok{x }\SpecialCharTok{{-}}\NormalTok{ ((redundant}\SpecialCharTok{*}\DecValTok{40}\NormalTok{)}\SpecialCharTok{/}\NormalTok{variant) }\SpecialCharTok{*}\NormalTok{ x }\SpecialCharTok{\^{}} \DecValTok{2} \SpecialCharTok{+} \FunctionTok{rnorm}\NormalTok{ (}\DecValTok{100}\NormalTok{)}
\end{Highlighting}
\end{Shaded}

\begin{Shaded}
\begin{Highlighting}[]
\NormalTok{loocv.fn }\OtherTok{=} \ControlFlowTok{function}\NormalTok{(df, exponent) \{}
\NormalTok{  glm.fit }\OtherTok{=} \FunctionTok{glm}\NormalTok{(y}\SpecialCharTok{\textasciitilde{}}\FunctionTok{poly}\NormalTok{(x, exponent), df)}
\NormalTok{\}}
\end{Highlighting}
\end{Shaded}

\begin{Shaded}
\begin{Highlighting}[]
\FunctionTok{set.seed}\NormalTok{(variant)}
\NormalTok{df }\OtherTok{\textless{}{-}} \FunctionTok{data.frame}\NormalTok{(x,y)}
\NormalTok{cv.error }\OtherTok{\textless{}{-}} \FunctionTok{rep}\NormalTok{(}\DecValTok{0}\NormalTok{, }\DecValTok{4}\NormalTok{)}

\ControlFlowTok{for}\NormalTok{ (i }\ControlFlowTok{in} \DecValTok{1}\SpecialCharTok{:}\DecValTok{4}\NormalTok{) \{}
\NormalTok{  glm.fit }\OtherTok{\textless{}{-}} \FunctionTok{glm}\NormalTok{(y}\SpecialCharTok{\textasciitilde{}}\FunctionTok{poly}\NormalTok{(x, i))}
\NormalTok{  cv.error[i] }\OtherTok{\textless{}{-}} \FunctionTok{cv.glm}\NormalTok{(df, glm.fit)}\SpecialCharTok{$}\NormalTok{delta[}\DecValTok{1}\NormalTok{]}
\NormalTok{\}}
\end{Highlighting}
\end{Shaded}

\begin{Shaded}
\begin{Highlighting}[]
\NormalTok{cv.error}
\end{Highlighting}
\end{Shaded}

\begin{verbatim}
## [1] 302.647785   1.186369   1.207521   1.246523
\end{verbatim}

\begin{Shaded}
\begin{Highlighting}[]
\FunctionTok{plot}\NormalTok{(}\DecValTok{1}\SpecialCharTok{:}\DecValTok{4}\NormalTok{, cv.error)}
\end{Highlighting}
\end{Shaded}

\includegraphics{lab4_files/figure-latex/unnamed-chunk-22-1.pdf}
Найменша помилка у квадратичної моделі. Це відповідає очікуванням,
оскільки у даних квадратична залежність.

\end{document}
